\documentclass[addpoints]{exam}
% \documentclass[addpoints,12pt]{exam}

\usepackage[slovene]{babel}
\usepackage{amsfonts}
\usepackage{amssymb}
\usepackage{amsmath}
\usepackage{float}
\usepackage{graphicx}
\usepackage{enumitem}
\usepackage{color}
\usepackage{eurosym}


%Komentar naslednje vrstice glede na verzijo testa -- rešitve ali prostor za odgovore:
% \printanswers

% \linespread{2}


\pointsinrightmargin
\marginbonuspointname{}


\renewcommand{\solutiontitle}{\noindent\textbf{Rešitve in točkovanje:}\par\noindent}

\colorgrids
\definecolor{GridColor}{gray}{0.7}
\setlength{\gridsize}{7mm}

\extrawidth{5mm}
\setlength{\rightpointsmargin}{1.5cm}

\begin{document}

\runningfooter{}{}{}

\begin{center}
    \fbox{\fbox{\parbox{15cm}{\centering
    Tretje preverjanje znanja \\ 1. letnik \\ februar 2026 \\ (Racionalna števila)}}}
\end{center}

\vspace{0.1in}
\makebox[\textwidth]{Ime in priimek, oddelek:\enspace\hrulefill}


\begin{center}
    \chqword{Naloga}
    \chpword{Možne točke}
    \chbpword{Dodatne točke}
    \chsword{Dosežene točke}
    \chtword{Skupaj}  
    \combinedpointtable[h][questions]
    % \combinedgradetable[h][questions]

\end{center}

\vspace{0.1in}


\begin{questions}

    % 5. vprašanje
    \question
    Izračunajte vrednost izraza.
    \begin{parts}
        \part[5]
        $\frac{9}{5}\cdot \frac{10}{12}-\dfrac{3\frac{2}{8}+\frac{5}{4}}{\frac{22}{12}}$

        % \begin{solution}
        %     \begin{align*} 
        %         \frac{18}{5}\cdot 2\frac{1}{12}-\dfrac{\frac{27}{8}+1\frac{5}{6}}{\frac{19}{8}}&=\frac{18}{5}\cdot \frac{25}{12}-\dfrac{\frac{27}{8}+\frac{11}{6}}{\frac{19}{8}}= \\
        %                                                                                     &= \frac{3}{1}\cdot \frac{5}{2}-\dfrac{\frac{108}{24}+\frac{44}{24}}{\frac{19}{8}}= \\
        %                                                                                     &=\frac{15}{2}-\dfrac{\frac{152}{24}}{\frac{19}{8}}= \\
        %                                                                                     &=\frac{15}{2}-\frac{152\cdot 8}{24\cdot 19}= \\
        %                                                                                     &=\frac{15}{2}-\frac{8}{3}= \\
        %                                                                                     &=\frac{45}{6}-\frac{16}{6}= \\
        %                                                                                     &=\frac{29}{6}
        %     \end{align*}
        %     Za pravilno množenje ulomkov $1$ točka, pravilno seštevanje ulomkov $1$ točka, pravilno razreševanje dvojnih ulomkov $1$ točka, pravilno odštevanje ulomkov $1$ točka, pravilen okrajšan rezultat $1$ točka.
        % \end{solution}


        \part[5]
        $1.\overline{5}:(0.5+0.\overline{4})-0.3\overline{22}\cdot 22$

        % \begin{solution}
        %     \begin{align*}
        %         2.\overline{6}:(0.5-0.\overline{3})-0.3\overline{72}\cdot 55 &= 2\frac{2}{3}:(\frac{1}{2}-\frac{1}{3})-\frac{41}{110}\cdot 55= \\
        %                                                                     &= \frac{8}{3}:\frac{1}{6}-\frac{41}{2}= \\
        %                                                                     &= \frac{8}{3}\cdot \frac{6}{1}-\frac{41}{2}= \\
        %                                                                     &= 16-\frac{41}{2}= \\
        %                                                                     &= -\frac{9}{2}
        %     \end{align*}
        %     Za pravilno zapisano število $0.3\overline{72}$ z ulomkom $1$ točka, pravilno zapisani števili $2.\overline{6}$ in $0.\overline{3}$ z ulomkoma $1$ točka, pravilno množenje in deljenje ulomkov $1$ točka, pravilno odštevanje ulomkov $1$ točka, končen rezultat $1$ točka.
        % \end{solution}

    \end{parts}



    % \newpage
    % 6. vprašanje

    \question
    Okrajšajte ulomke.
    \begin{parts}
        

        \part[4]
        $(x^2-4)^{-1}(x^2+5x+6)$

        % \begin{solution}
        %     \begin{align*}
        %         5(81x^2-1)(45x-5)^{-1}&= \dfrac{5(81x^2-1)}{45x-5}= \\
        %                             &= \dfrac{5(9x-1)(9x+1)}{5(9x-1)}= \\
        %                             &= 9x+1
        %     \end{align*}
        %     Pravilen zapis z ulomkom $1$ točka, razstavljanje razlike kvadratov in izpostavljanje skupnega faktorja $1$ točka, krajšanje ulomka $1$ točka, končen rezultat $1$ točka.
        % \end{solution}


        \part[5]
        $\dfrac{a^{n+2}-a^{n}}{a^{n-1}+a^n}$

        % \begin{solution}
        %     \begin{align*}
        %         \dfrac{x^{n+1}+2x^n-3x^{n-1}}{x^n-9x^{n-2}} &= \dfrac{x^{n-1}(x^{2}+2x-3)}{x^{n-2}(x^2-9)}= \\
        %                                                     &= \dfrac{x^{n-1}(x+3)(x-1)}{x^{n-2}(x-3)(x+3)}= \\
        %                                                     &= \dfrac{x(x-1)}{x+3}
        %     \end{align*}
        %     Za pravilno izpostavljanje skupnega faktorja v števcu in imenovalcu po $1$ točka, razstavljanje po Vietu in razlike kvadratov po $1$ točka, zapis pravilno okrajšanega ulomka $1$ točka.
        % \end{solution}


        \part[5]
        $\dfrac{(\frac{1}{5})^{-1}x^3\cdot 8^{-2}(x^{-2})^5}{(5x^{-2}8^{-2})^3}; \quad a,b\neq 0$

        % \begin{solution}
        %     \begin{align*}
        %         \dfrac{(\frac{1}{3})^{-1}a^2\cdot 10^{-3}(a^{-2})^4 (10^{-4})^3 (a^0+11b^0)}{(6a^{-3}10^{-8})^2} &= \dfrac{3a^2\cdot 10^{-3}a^{-8} 10^{-12} (1+11\cdot 1)}{6^2a^{-6}10^{-16}}= \\
        %                                                                                                         &= \dfrac{3\cdot 12 a^{2-8}10^{-3-12}}{36a^{-6}10^{-16}}= \\
        %                                                                                                         &= a^{-6-(-6)}10^{-15-(-16)}= \\
        %                                                                                                         &= a^0 10= \\
        %                                                                                                         &= 10
        %     \end{align*}
        %     Za uporabo pravila potenciranja potence $1$ točka, pravilno obratna vrednost ulomka $1$ točka, uporaba pravila množenja potenc z enakimi osnovami $1$ točka, pravilno upoštevanje $x^0=1$ $1$ točka, končen rezultat $1$ točka.
        % \end{solution}


    \end{parts}



    % \newpage

        % 7. vprašanje

        \question
        Poenostavite izraze.
        \begin{parts}
            

            \part[4]
            $\dfrac{x^3-25}{x^2+3x+2}-\left(\left(-\dfrac{x+1}{26}\right)^{-1}+\dfrac{25}{x+2}\right)$
    
            % \begin{solution}
            %     \begin{align*}
            %         \dfrac{x+4}{x+5}-\dfrac{2x-10}{x^2-x-12}:\dfrac{x^2-25}{x-4} &= \dfrac{x+4}{x+5}-\dfrac{2(x-5)}{(x-4)(x+3)}\cdot\dfrac{x-4}{(x-5)(x+5)}= \\
            %                                                                     &= \dfrac{x+4}{x+5}-\dfrac{2}{(x+3)(x-5)}= \\
            %                                                                     &= \dfrac{(x+4)(x+3)-2}{(x+5)(x+3)}= \\
            %                                                                     &= \dfrac{x^2+7x+10}{(x+5)(x+3)}= \\
            %                                                                     &= \dfrac{(x+5)(x+2)}{(x+5)(x+3)}= \\
            %                                                                     &=\dfrac{x+2}{x+3}
            %     \end{align*}
            %     Za pravilno faktorizacijo izrazov $1$ točka, pravilno deljenje ulomkov $1$ točka, pravilno odštevanje ulomkov $1$ točka, pravilno zapisan okrajšan rezultat $1$ točka.
            % \end{solution}
    
    
            \part[5]
            $\dfrac{\dfrac{x-2}{x^2-1}}{\dfrac{2-x}{1-x}}-(x+1)^{-1}$
    
            % \begin{solution}
            %     \begin{align*}
            %         \dfrac{(a-1)^{-1}-(a+1)^{-1}}{a+\frac{a}{a^2-1}} &= \dfrac{\frac{1}{a-1}-\frac{1}{a+1}}{\frac{a(a^2-1)+a}{a^2-1}}= \\
            %                                                         &= \dfrac{\frac{a+1-(a-1)}{a^2-1}}{\frac{a^3-a+a}{a^2-1}}= \\
            %                                                         &= \dfrac{\frac{2}{a^2-1}}{\frac{a^3}{a^2-1}}= \\
            %                                                         &= \dfrac{2(a^2-1)}{a^3(a^2-1)}= \\
            %                                                         &= \dfrac{2}{a^3}
            %     \end{align*}
            %     Za pravilen zapis obratnih vrednosti ulomkov $1$ točka, pravilno seštevanje in odštevanje ulomkov $1$ točka, pravilno razrešen dvojni ulomek $1$ točka, pravilno krajšanje ulomkov $1$ točka, končen rezultat $1$ točka.
            % \end{solution}
    
    
            \part[5]
            $\dfrac{x^{-2}+x^{-1}}{x^{-2}-x^{-1}}-(1-x)^{-1}$
    
            % \begin{solution}[3cm]
            %     \begin{align*}
            %         \dfrac{2x^{-2}}{1-2x^{-1}+x^{-2}}+\dfrac{x^{-1}}{1-x^{-1}} &= \dfrac{\frac{2}{x^2}}{1-\frac{2}{x}+\frac{1}{x^2}}+\dfrac{\frac{1}{x}}{1-\frac{1}{x}}= \\
            %                                                                 &= \dfrac{\frac{2}{x^2}}{\frac{x^2-2x+1}{x^2}}+\dfrac{\frac{1}{x}}{\frac{x-1}{x}}= \\
            %                                                                 &= \dfrac{2x^2}{(x^2-2x+1)x^2}+\dfrac{x}{(x-1)x}= \\
            %                                                                 &= \dfrac{2}{(x-1)^2}+\dfrac{1}{x-1}= \\
            %                                                                 &= \dfrac{2+(x-1)}{(x-1)^2}= \\
            %                                                                 &= \dfrac{x+1}{(x-1)^2}
            %     \end{align*}
            %     Za pravilen zapis obratnih vrednosti $1$ točka, pravilno seštevanje in odštevanje ulomkov $1$ točka, pravilno razreševanje dvojnih ulomkov $1$ točka, pravilno krajšanje ulomkov $1$ točka, končen rezultat $1$ točka.
            % \end{solution}
    
    
        \end{parts}




        \question[4]
    Okrajšajte ulomek: $\dfrac{b^{2x+y}+b^{x+2y}+b^{x+y+2}}{b^{2x-y}+b^x+b^{x-y+2}}$.

    % \begin{solution}[7cm]
    %     \begin{align*}
    %         \dfrac{b^{2x+y}+b^{x+2y}+b^{x+y+2}}{b^{2x-y}+b^x+b^{x-y+2}} &= \dfrac{b^{x+y}(b^x+b^y+b^2)}{b^{x-y}(b^x+b^y+b^2)}= \\
    %                                                                     &= \dfrac{b^{x+y}}{b^{x-y}}= \\
    %                                                                     &= b^{2y}
    %     \end{align*}
    %     Za pravilno izpostavljanje skupnega faktorja $1$ točka; za pravilno krajšanje $1$ točka; končen rezultat $1$ točka.
    % \end{solution}




        % 9. vprašanje
    \question[4]
    V posodi imamo $220~g$ $95~\%$ raztopine kisline. Koliko gramov $60~\%$ raztopine kisline moramo dodati, da bomo dobili $80~\%$ raztopino kisline?
    
    
    % \begin{solution}[7cm]
    %     ???
    % \end{solution}



    % \newpage

    % 10. vprašanje
    \question[4]
    Izlet v Grčijo se je najprej podražil za $5~\%$, nato pa še za $10~\%$. 
    Koliko je stal na začetku, če zdaj stane $2310~€$?
    
    
    % \begin{solution}[3cm]
    %     ???
    % \end{solution}


\end{questions}



\end{document}